% Generated draft: requires mathematical and source review.
\documentclass[11pt]{article}
\usepackage[margin=25mm]{geometry}
\usepackage{fontspec}
\setmainfont{Latin Modern Roman}
\usepackage{amsmath,amssymb,mathtools,unicode-math}
\setmathfont{Latin Modern Math}
\usepackage{enumitem,needspace,xurl,hyperref}
\setlength{\parindent}{0pt}
\setlength{\parskip}{5pt}
\setlength{\emergencystretch}{3em}
\newcommand{\N}{\mathbb{N}}
\newcommand{\Z}{\mathbb{Z}}
\newcommand{\Q}{\mathbb{Q}}
\newcommand{\R}{\mathbb{R}}
\newcommand{\C}{\mathbb{C}}
\newcommand{\F}{\mathbb{F}}
\newcommand{\A}{\mathbb{A}}
\newcommand{\PP}{\mathbb{P}}
\newcommand{\E}{\mathbb{E}}
\newcommand{\eps}{\varepsilon}
\newcommand{\dd}{\,\mathrm{d}}
\begin{document}
\section*{UnsolvedMath ID: 10800006}
\textbf{Source problem code:} catalogue entry
\subsection{Definitions and mathematical statement}
For a real smooth function on an $n$-manifold, the Morse index of a nondegenerate critical point is the number of negative eigenvalues of its Hessian. The supplied registry explains the otherwise unspecified phrase ``these indices'' as the integer Morse indices of two critical points created in a Morse birth surgery.

In that intended setting, let $f_t$ be the relevant one-parameter deformation, with no newborn pair at $t=-\varepsilon$ and the pair $p_+,q_+$ at $t=+\varepsilon$, for sufficiently small $\varepsilon>0$. At a generic birth their indices are adjacent integers $k,k+1$. Vassiliev's Problem 1F asks whether a convenient topological characteristic of the pre-birth function $f_{-\varepsilon}$ can predict these integers, beyond the parity information reported as already available.

An eventual answer would specify allowable topological data $T(f_{-\varepsilon})$ and a rule recovering the unordered pair
\[
 \{\operatorname{ind}_{p_+}(f_{+\varepsilon}),
   \operatorname{ind}_{q_+}(f_{+\varepsilon})\}
\]
for the admissible birth surgeries. The excerpt does not give the ambient dimension, function/deformation class, local domain and boundary data, the birth location or other information retained, or the meaning of topological equivalence of pre-birth functions. Accordingly no assertion is made that an unspecified global topological type alone determines the indices, nor is a particular characteristic substituted for the requested one.
\subsection{Short English statement}
Can topological data available before a Morse birth predict the exact integer indices of the two newborn critical points, rather than only their parity? The permitted data and family hypotheses are missing from the excerpt.
\subsection{Sources}
\begin{itemize}
\item Local registry: \url{data/unsolvedmath/problems.json}, permanent UnsolvedMath ID $10800006$, exact source code \texttt{catalogue entry}, attributed to Victor Vassiliev.
\item Vassiliev, \emph{A Few Problems on Monodromy and Discriminants}, Problem 1F: \url{https://amj.math.stonybrook.edu/html-articles/Files-2015-2024/15-11/}. Supplied access date: 2026-07-29; extraction: html-structure; eighteen labeled problem environments.
\item Supplied background reference: Vassiliev, \emph{Applied Picard--Lefschetz Theory}, AMS, 2002, for the parity prediction in real Morsification theory. Supplied computational-work identifiers: \url{https://arxiv.org/abs/2109.12287} and \url{https://arxiv.org/abs/2512.12738}; titles and exact relevant results are not supplied.
\end{itemize}
\end{document}
